% Part: propositional-logic
% Chapter: syntax-and-semantics
% Section: formulas

\documentclass[../../../include/open-logic-section]{subfiles}

\begin{document}

\olfileid{pl}{syn}{fml}

\olsection{বচনমূলক \usetoken{P}{formula}}

বচনমূলক যুক্তিবিদ্যার !!^{formula}s গঠিত হয়
\emph{!!{propositional
    variable}s}\iftag{prvFalse}{\iftag{prvTrue}{,}{ ও}
  মিথ্যার বচনধ্রুবক~$\lfalse$}{}\iftag{prvTrue}{\iftag{prvFalse}{
    ও}{} সত্যের বচনধ্রুবক~$\ltrue$}{} থেকে, \emph{যৌক্তিক
  সংযোজক} ব্যবহার করে।

\begin{enumerate}
\item !!^a{denumerable} সেট~$\PVar$, যার !!{propositional variable}s হলো $\Obj p_0$,
  $\Obj p_1$, \dots
\tagitem{prvFalse}{!!{falsity}-র বচনধ্রুবক~$\lfalse$।}{}
\tagitem{prvTrue}{!!{truth}-এর বচনধ্রুবক~$\ltrue$।}{}
\item যৌক্তিক সংযোজকগুলি:
  \startycommalist
  \iftag{prvNot}{\ycomma $\lnot$ (নঞর্থকরণ)}{}%
  \iftag{prvAnd}{\ycomma $\land$ (সংযোজন)}{}%
  \iftag{prvOr}{\ycomma $\lor$ (বিয়োজন)}{}%
  \iftag{prvIf}{\ycomma $\lif$ (!!{conditional})}{}%
  \iftag{prvIff}{\ycomma $\liff$ (!!{biconditional})}{}%
\item যতি-সংকেত: (, ), এবং কমা।
\end{enumerate}

বচনমূলক যুক্তিবিদ্যার এই ভাষাটিকে আমরা $\Lang L_0$ দিয়ে নির্দেশ করি।

\iftag{defNot,defOr,defAnd,defIf,defIff,defTrue,defFalse,defEx,defAll}{%
উপরে প্রবর্তিত মৌলিক সংযোজকগুলির পাশাপাশি আমরা নিচের
\emph{সংজ্ঞায়িত} সংকেতগুলিও ব্যবহার করি:
\startycommalist
  \iftag{defNot}{\ycomma $\lnot$ (নঞর্থকরণ)}{}%
  \iftag{defAnd}{\ycomma $\land$ (সংযোজন)}{}%
  \iftag{defOr}{\ycomma $\lor$ (বিয়োজন)}{}%
  \iftag{defIf}{\ycomma $\lif$ (!!{conditional})}{}%
  \iftag{defIff}{\ycomma $\liff$ (!!{biconditional})}{}%
  \iftag{defFalse}{\ycomma $\lfalse$ (!!{falsity})}{}%
  \iftag{defTrue}{\ycomma $\ltrue$ (!!{truth})}}{}।

\begin{tagblock}{defNot,defOr,defAnd,defIf,defIff,defTrue,defFalse,defEx,defAll}
\begin{explain}
সংজ্ঞায়িত সংকেত সরকারি ভাবে ভাষাটির অংশ নয়; বরং অনানুষ্ঠানিক সংক্ষিপ্ত
রূপ হিসেবে সেটি প্রবর্তিত হয়। কেবল মৌলিক সংকেত ব্যবহার করলে যে সূত্রগুলি
বেশ দীর্ঘ হতো, এই সংকেত দিয়ে সেগুলিকে সংক্ষেপে লেখা যায়। এটি নিঃসন্দেহে
সুবিধাজনক। তবে আরও বড় সুবিধা হলো, প্রমাণগুলি ছোট হয়। কোনো সংকেত মৌলিক
হলে প্রমাণে সেটিকে আলাদাভাবে বিচার করতে হয়। তাই মৌলিক সংকেত যত বেশি,
আমাদের প্রমাণও তত দীর্ঘ হয়।
\end{explain}
\end{tagblock}

% Alternate symbols

\begin{intro}
উপরে ব্যবহৃত পরিভাষা ও সংকেতগুলির বদলে অন্য পরিভাষা ও সংকেতের সঙ্গে তোমার
পরিচয় থাকতে পারে। যুক্তিবিদ্যার বইগুলি (এবং শিক্ষকরা) “নঞর্থকরণ”-এর জন্য
সাধারণত $\sim$, $\neg$ ও~!\; “সংযোজন”-এর জন্য $\wedge$, $\cdot$ ও $\&$
ব্যবহার করেন। “শর্তবচন” বা “নিহিতকরণ”-এর প্রচলিত সংকেত হলো
$\rightarrow$, $\Rightarrow$ ও $\supset$।
\iftag{prvIff,defIff}{“দ্বিশর্তবচন”, “দ্বিমুখী নিহিতকরণ” বা “(বস্তুগত)
  সমতুল্যতা”-র সংকেত হলো $\leftrightarrow$,
  $\Leftrightarrow$ ও $\equiv$।}{}
\iftag{prvFalse,defFalse}{$\lfalse$ সংকেতটিকে নানা ক্ষেত্রে
  “মিথ্যা”, “মিথ্যাধ্রুবক”, “অসঙ্গতি” বা “নিম্ন” বলা হয়।}{}
\iftag{prvTrue,defTrue}{$\ltrue$ সংকেতটিকে নানা ক্ষেত্রে
  “সত্য”, “সত্যধ্রুবক” বা “ঊর্ধ্ব” বলা হয়।}{}
\end{intro}

\begin{defn}[সূত্র]
\ollabel{defn:formulas}
বচনমূলক যুক্তিবিদ্যার \emph{!!{formula}s}-এর সেট~$\Frm[L_0]$
নিচের মতো আরোহীভাবে সংজ্ঞায়িত:
\begin{enumerate}
\tagitem{prvFalse}{$\lfalse$ একটি পরমাণু !!{formula}।}{}

\tagitem{prvTrue}{$\ltrue$ একটি পরমাণু !!{formula}।}{}

\item প্রতিটি !!{propositional variable}~$\Obj p_i$ একটি পরমাণু
  !!{formula}।

\tagitem{prvNot}{যদি $!A$ !!a{formula} হয়, তবে $\lnot !A$
  !!a{formula}।}{}

\tagitem{prvAnd}{যদি $!A$ ও $!B$ !!{formula}s হয়, তবে $(!A \land
  !B)$ !!a{formula}।}{}

\tagitem{prvOr}{যদি $!A$ ও $!B$ !!{formula}s হয়, তবে $(!A \lor !B)$
  !!a{formula}।}{}

\tagitem{prvIf}{যদি $!A$ ও $!B$ !!{formula}s হয়, তবে $(!A \lif !B)$
  !!a{formula}।}{}

\tagitem{prvIff}{যদি $!A$ ও $!B$ !!{formula}s হয়, তবে $(!A \liff !B)$
  !!a{formula}।}{}

\tagitem{limitClause}{অন্য কিছুই !!a{formula} নয়।}{}
\end{enumerate}
\end{defn}

\begin{explain}
!!{formula}s-এর সংজ্ঞাটি একটি
\emph{আরোহী সংজ্ঞা}। মূলত অসীমসংখ্যক ধাপে আমরা
!!{formula}s-এর সেটটি নির্মাণ করি। প্রথম ধাপে সব পরমাণু সূত্রকে সূত্র
ঘোষণা করি; সংজ্ঞার প্রথম কয়েকটি বিধি, অর্থাৎ
\iftag{prvTrue}{$\ltrue$, }{}%
\iftag{prvFalse}{$\lfalse$, }{}%
$\Obj p_i$-এর বিধিগুলি এই ধাপের সঙ্গে মেলে। তাই “পরমাণু !!{formula}”
বলতে এই রূপের যেকোনো !!{formula} বোঝায়।

সংজ্ঞার অন্য বিধিগুলি আগে নির্মিত !!{formula}s থেকে নতুন
!!{formula}s নির্মাণের নিয়ম দেয়। দ্বিতীয় ধাপে নিয়মগুলি ব্যবহার করে
পরমাণু !!{formula}s থেকে !!{formula}s নির্মাণ করা যায়। তৃতীয় ধাপে
পরমাণু সূত্র এবং দ্বিতীয় ধাপে পাওয়া সূত্রগুলি থেকে নতুন সূত্র নির্মাণ করি;
তারপর একইভাবে এগিয়ে যাই। কোনো ধাপে শেষ পর্যন্ত যা নির্মিত হয়, এবং শুধু
তাই, !!{formula}।
\end{explain}

$!B$, $!C$ থেকে দ্বিস্থানী সংযোজক~$\ast$ ব্যবহার করে নির্মিত
$(!B \ast !C)$ সূত্রটি লেখার সময় আমরা প্রায়ই বাইরের এক জোড়া বন্ধনী
বাদ দিয়ে শুধু~$!B \ast !C$ লিখব।

\begin{tagblock}{defNot,defOr,defAnd,defIf,defIff,defTrue,defFalse,defEx,defAll}
\begin{defn}
সংজ্ঞায়িত কারকগুলি ব্যবহার করে নির্মিত সূত্রগুলিকে নিচের মতো বুঝতে হবে:

\begin{tagenumerate}{defTrue,defFalse,defNot,defOr,defAnd,defIf,defIff,defEx,defAll}

\tagitem{defTrue}{$\ltrue$ হলো
  \iftag{prvFalse}{$\lnot\lfalse$}{কোনো নির্দিষ্ট পরমাণু
    !!{formula}~$!A$-এর জন্য $(!A \lor \lnot !A)$}-এর সংক্ষিপ্ত রূপ।}{}

\tagitem{defFalse}{$\lfalse$ হলো
  \iftag{prvTrue}{$\lnot\ltrue$}{কোনো নির্দিষ্ট পরমাণু
    !!{formula}~$!A$-এর জন্য $(!A \land \lnot !A)$}-এর সংক্ষিপ্ত রূপ।}{}

\tagitem{defNot}{$\lnot !A$ হলো $!A \lif \lfalse$-এর সংক্ষিপ্ত রূপ।}{}

\tagitem{defOr}{$!A \lor !B$ হলো
  \iftag{prvAnd}{$\lnot(\lnot !A \land \lnot !B)$}{$\lnot !A \lif
    !B$}-এর সংক্ষিপ্ত রূপ।}{}

\tagitem{defAnd}{$!A \land !B$ হলো
  \iftag{prvOr}{$\lnot(\lnot !A \lor \lnot !B)$}{$\lnot (!A \lif
    \lnot !B)$}-এর সংক্ষিপ্ত রূপ।}{}

\tagitem{defIf}{$!A \lif !B$ হলো
  \iftag{prvOr}{$(\lnot !A \lor !B)$}{$\lnot (!A \land \lnot !B)$}-এর সংক্ষিপ্ত রূপ।}{}
% BN-SRC-019 TARGET-CORRECTION-BEGIN
\par\smallskip\noindent\textnormal{\small\textbf{উৎস-সংশোধন (BN-SRC-019):} হিমায়িত ইংরেজি উৎসে প্রথম বিকল্পটির সূত্র $\lnot !A \lor !B)$ লেখা, যেখানে শুধু সমাপনী বন্ধনী আছে। গঠন-বিধি ও সংজ্ঞাটির অন্য বিকল্প অনুসারে বাংলা পাঠে সুষম সূত্র $(\lnot !A \lor !B)$ ব্যবহার করা হয়েছে।} % BN-SRC-019 TARGET-CORRECTION-END

\tagitem{defIff}{$!A \liff !B$ হলো $(!A \lif !B) \land (!B
  \lif !A)$-এর সংক্ষিপ্ত রূপ।}{}
\end{tagenumerate}
\end{defn}
\end{tagblock}

\begin{defn}[সংকেতবিন্যাসগত অভিন্নতা]
$\ident$ সংকেতটি প্রতীকক্রমগুলির মধ্যে সংকেতবিন্যাসগত অভিন্নতা প্রকাশ করে;
অর্থাৎ $!A \ident !B$ তখনই, যখন $!A$ ও $!B$ সমান দৈর্ঘ্যের প্রতীকক্রম এবং
প্রতিটি স্থানে তাদের একই সংকেত থাকে।
\end{defn}

সংযুক্তকরণের ফলে পাওয়া প্রতীকক্রমও $\ident$ সংকেতটির দুপাশে থাকতে পারে।
যেমন, $!A \ident (!B \lor !C)$-এর অর্থ: $!A$ প্রতীকক্রমটি সেই প্রতীকক্রমের
সঙ্গে অভিন্ন, যা এই ক্রমে একটি প্রারম্ভিক বন্ধনী, $!B$ প্রতীকক্রম,
$\lor$ সংকেত, $!C$ প্রতীকক্রম এবং একটি সমাপনী বন্ধনী সংযুক্ত করে পাওয়া
যায়। যদি তা-ই হয়, তবে আমরা জানি যে $!A$-এর প্রথম সংকেত একটি প্রারম্ভিক
বন্ধনী; $!A$-তে $!B$ একটি উপপ্রতীকক্রম হিসেবে আছে (দ্বিতীয় সংকেত থেকে
শুরু), তার পরে $\lor$ আছে, ইত্যাদি।

\end{document}
